\section{Propiedad intelectual}

El proyecto se acoge al Reglamento de Propiedad Intelectual de EAFIT, aprobado mediante acta 002 del 10 de abril de 2025 \cite{eafit_pi,eafit_pi_web}. Los \emph{derechos morales} son inalienables e irrenunciables y corresponden a quien realiza la creación: Alexander Saldarriaga Vélez como autor, y el asesor en la medida de su aporte intelectual (art.~6, par.~1; art.~20, num.~1).

Sobre los \emph{derechos patrimoniales}, el artículo~7 fija como regla general que los estudiantes son titulares de lo que realicen en actividades académicas, \emph{incluidos los trabajos de grado}, salvo contribuciones significativas de EAFIT o de terceros. Dos disposiciones restringen esa regla aquí: el artículo~6, numeral~2 atribuye la titularidad a EAFIT cuando se usan \emph{recursos significativos} —que el artículo~5, numeral~12 define e incluyen laboratorios, talleres y equipos—, y el proyecto usará el taller y el Laboratorio de Mecatrónica; y el parágrafo del artículo~7 asigna a la Universidad la fracción correspondiente al aporte intelectual del asesor. Se prevé, en consecuencia, titularidad patrimonial compartida. Como ese parágrafo exige que cualquier reparto distinto se acuerde por escrito \emph{antes del inicio de las actividades}, el estudiante y el asesor radicarán el porcentaje de cada parte (art.~19, num.~11) ante la jefatura de Transferencia de Tecnología y Conocimiento antes de comenzar la ejecución.

La protección aplicable al documento, los planos, el modelo CAD y el firmware es la de derechos de autor, sin registro constitutivo; la decisión de registrar corresponde a EAFIT (art.~10). No se prevé patente: la novedad está en la configuración y el ajuste experimental de tecnologías conocidas. Si apareciera materia patentable, se validará con esa área antes de divulgarla (arts.~11 y~12). El prototipo se transfiere al Laboratorio de Mecatrónica para uso docente e interno, sin contraprestación.
